% =====================================================================
% DIAL — Supplementary Material
% Target venue: Bioinformatics (OUP), supplementary appendix.
% =====================================================================
\documentclass[11pt,a4paper]{article}
\usepackage[margin=1in]{geometry}
\usepackage{booktabs,amsmath,amssymb,amsthm}
\usepackage{graphicx}
\usepackage{hyperref}
\usepackage{natbib}
\usepackage{xcolor}
\usepackage{array}

\renewcommand{\thesection}{S\arabic{section}}
\renewcommand{\thetable}{S\arabic{table}}
\renewcommand{\thefigure}{S\arabic{figure}}

\newtheorem{lemma}{Lemma}
\newtheorem{corollary}{Corollary}

\title{DIAL --- Supplementary Material \\
\large \textit{a direction-invariant audit metric for batch-corrected
cancer subtype classifiers}}
\author{Seungho Kuk}
\date{2026-04-25}

\begin{document}
\maketitle
\tableofcontents
\vspace{1em}\hrule\vspace{1em}

\paragraph{Overview.} This supplementary document collects seven
reviewer-facing robustness analyses commissioned for the v5.1 DIAL
paper, a complete proof of the label-flip lemma, and a
code/data-availability statement. Each section asks: if we change the
batch-correction estimator (S2), the DE test (S5), the classifier
paradigm (S6), or the correction family (S7), does the thyroid-specific
DIAL flip survive? Bottom line: the THCA flip is robust under the
count-native ComBat-seq estimator (S2), is the only cohort flagged by
random-effects meta-analysis (S3), is a gene-level phenomenon that
pathway aggregation eliminates (S4), is not an artefact of the Welch
$t$-test used upstream (S5), persists under a variational quantum
classifier (S6), and reshuffles sensibly under FastRNA, BUHMBOX, and
LMM Han-lab methods (S7).

% =====================================================================
\section{Mutation label schemes and exclusions}
\label{sec:S1}

The cohort availability matrix is given in Table~1 of the main text;
here we document only the mutation-label conventions and the
exclusion reasons for cohorts not carried through to the DIAL
analysis.

\paragraph{Subtype labels by cancer.}
\begin{itemize}
\item \textbf{THCA} --- BRAF V600E vs RAS (NRAS/HRAS/KRAS) hotspot
  mutations. Wild-type and RET/PTC fusion samples are excluded.
\item \textbf{SKCM} --- BRAF vs NRAS hotspot mutations (mutually
  exclusive in TCGA-SKCM).
\item \textbf{LGG} --- IDH1/IDH2 mutation status (mutated vs
  wild-type).
\item \textbf{LUAD} --- KRAS vs EGFR mutation status (mutually
  exclusive when both are called in the same sample, resolved to the
  higher-VAF call).
\item \textbf{COAD} --- MSI-H vs MSS (mismatch-repair status from
  TCGA MC3 or GEO companion annotation).
\end{itemize}

\paragraph{Exclusion reasons.} Seven cohorts considered during Phase
1 were excluded: \texttt{GSE33630, GSE29265} (THCA, no mutation
labels); \texttt{GSE65904} (SKCM, no labels); \texttt{GSE16011} (LGG,
no labels); \texttt{GSE4271} (LGG, HTTP 404 on GEO FTP);
\texttt{GSE72094} (LUAD, GPL15048 probe--gene mapping failed);
\texttt{GSE17536} (COAD, no MSI/MSS labels).

% =====================================================================
\section{ComBat-seq robustness}
\label{sec:S2}

We re-ran the v5.1 LODO DIAL pipeline on THCA with two
batch-correction variants: (a) the v5.1 baseline, ComBat
(\texttt{pycombat\_norm}) on log2-TPM; and (b) ComBat-seq
\citep{zhang2020combatseq} via \texttt{inmoose.pycombat.pycombat\_seq}
on raw STAR unstranded counts for the RNA-seq cohort, followed by a
cross-platform ComBat step to bridge to the Affymetrix microarray
values. ComBat-seq is mathematically defined only on RNA-seq counts;
we therefore cannot run it jointly across a count cohort and a
microarray cohort, and the hybrid above is the strongest defensible
construction given that THCA has exactly one cohort of each type.

Four of five classifiers show $|\Delta\mathrm{DIAL}| < 0.1$:
LogReg-$\ell_2$ $0.494 \to 0.495$,
LogReg-elasticnet $0.492 \to 0.494$,
GradientBoosting $0.334 \to 0.273$,
XGBoost $0.013 \to 0.000$. Only RandomForest shrinks appreciably
($0.324 \to 0.129$; post-correction AUC $0.177 \to 0.372$), and its
AUC remains well below $0.5$. No classifier recovers BRAF/RAS signal
above chance under either correction. The flip is therefore not an
artefact of log2-TPM handling; switching to the count-native estimator
leaves four of five classifiers in the same DIAL regime. (Artefact:
\texttt{results/v8\_statgen/v8\_combatseq\_vs\_combat.tsv}.)

% =====================================================================
\section{Random-effects meta-analysis}
\label{sec:S3}

\paragraph{Background.} The primary analysis measured DIAL for five
classifiers on five TCGA-derived cohorts. To formalise that the
thyroid signal is specific rather than a heavy-tailed pan-cancer
phenomenon, we applied random-effects meta-analysis in the style of
METASOFT \citep{han2011metasoft,han2012mvalue,lee2017metasoft}.

\paragraph{Methods.} Per-classifier, per-cancer DIAL is the effect
size $\theta_i$ ($i = 1..5$). $\mathrm{SE}(\mathrm{DIAL})$ is
approximated by Hanley--McNeil $\mathrm{SE}(\mathrm{AUC}_{\mathrm{post}})$
using class counts from the harmonisation table; DIAL is a
shifted/flipped function of AUC$_{\mathrm{post}}$ so this is valid to
first order. SE is floored at $10^{-4}$ to avoid divide-by-zero under
perfect separation. Between-study $\tau^2$ uses DerSimonian--Laird;
Cochran's $Q$ is tested against $\chi^2$ with $k - 1 = 4$ d.f.;
$I^2 = \max(0, (Q - \mathrm{df})/Q)$. Han $m$-values
\citep{han2012mvalue} use a Gaussian two-component mixture with a
flat $0.5$ prior: $P(\mathrm{data} \mid \mathrm{effect})$ at the RE
pooled mean with variance $\mathrm{SE}_i^2 + \tau^2$ versus
$P(\mathrm{data} \mid \mathrm{no\ effect})$ at zero with variance
$\mathrm{SE}_i^2$.

\paragraph{Results.}
\begin{table}[htbp]
\centering
\small
\begin{tabular}{lrrrrrrrrrr}\toprule
Classifier & $\bar\theta$ & $\tau^2$ & $Q$ & $p(Q)$ & $I^2$ &
$m_{\mathrm{THCA}}$ & $m_{\mathrm{SKCM}}$ & $m_{\mathrm{LGG}}$ &
$m_{\mathrm{LUAD}}$ & $m_{\mathrm{COAD}}$ \\ \midrule
LogReg-$\ell_2$      & 0.099 & 0.114 & 4538.8 & $<10^{-16}$ & 99.9\% &
1.00 & 0.10 & 0.02 & 0.04 & 0.03 \\
RandomForest         & 0.065 & 0.017 & 184.1 & $<10^{-16}$ & 97.8\% &
1.00 & 0.22 & 0.06 & 0.16 & 0.10 \\
XGBoost              & 0.001 & 0.000 & 0.12  & 0.998       & 0.0\%  &
0.50 & 0.50 & 0.50 & 0.50 & 0.50 \\
LogReg-elasticnet    & 0.099 & 0.110 & 4471.9 & $<10^{-16}$ & 99.9\% &
1.00 & 0.10 & 0.02 & 0.04 & 0.04 \\
GradientBoosting     & 0.076 & 0.026 & 186.6 & $<10^{-16}$ & 97.9\% &
1.00 & 0.18 & 0.07 & 0.28 & 0.12 \\
\bottomrule\end{tabular}
\caption{Per-classifier random-effects meta-analysis of DIAL across
five cancers. Four of five classifiers reject homogeneity
overwhelmingly with THCA $m \geq 0.9999$; XGBoost is a null control.}
\label{tab:s3}
\end{table}

Four classifiers reject homogeneity with $I^2 > 97\%$ and
$p(Q) < 10^{-16}$. In every rejected case THCA has $m \geq 0.9999$
while every other cohort sits at or below the Han $m \leq 0.1$
no-effect band. XGBoost is a negative control: $Q$ is non-significant
($p = 0.998$), $I^2 = 0$, and every $m$-value collapses to the $0.5$
flat prior --- consistent with XGBoost never flipping in the primary
analysis.

% =====================================================================
\section{Pathway-level DIAL (MSigDB Hallmark)}
\label{sec:S4}

The v5.1 primary analysis measured DIAL on the $3{,}000$
variance-top genes, where THCA BRAF-vs-RAS produced a $0.494$ flip
post-ComBat. The same $392$ THCA samples, scored instead as a
50-dimensional MSigDB Hallmark \citep{liberzon2015hallmark} ssGSEA
matrix \citep{barbie2009ssgsea, hanzelmann2013gsva}, return pooled
$\mathrm{DIAL} = 0.000$ across all five classifiers with
$\mathrm{AUC}_{\mathrm{post}} \in [0.90, 0.98]$.

\begin{table}[htbp]
\centering
\small
\begin{tabular}{lrrrr}\toprule
Classifier & DIAL gene & DIAL pathway & AUC$_{\mathrm{post}}$ gene &
AUC$_{\mathrm{post}}$ pathway \\ \midrule
LogReg-$\ell_2$        & 0.494 & 0.000 & 0.006 & 0.964 \\
LogReg-elasticnet      & 0.492 & 0.000 & 0.008 & 0.958 \\
RandomForest           & 0.323 & 0.000 & 0.177 & 0.977 \\
GradientBoosting       & 0.334 & 0.000 & 0.166 & 0.918 \\
XGBoost                & 0.013 & 0.000 & 0.487 & 0.905 \\
\bottomrule\end{tabular}
\caption{Pooled gene-level versus pathway-level DIAL on THCA.}
\label{tab:s4a}
\end{table}

To distinguish aggregation-dilutes (some pathways flip but the pooled
vote cancels) from pathway-level resilience (no pathway flips), we
computed per-pathway LODO DIAL: one LogReg-$\ell_2$ per Hallmark
column, same joint 50-D ComBat. Zero of 50 pathways cross the
batch-entangled threshold; median DIAL $= 0.000$, maximum $= 0.085$
(Oxidative Phosphorylation, the canonical RAS-high signature in
thyroid carcinoma \citep{landa2016thyroid}). $29/50$ pathways are
clean \texttt{true\_biology}, $13/50$ \texttt{ambiguous}, $8/50$
\texttt{no\_signal}. The v5.1 $3{,}000$-gene DIAL $0.494$ is therefore
a high-dimensional residual-batch phenomenon, not a reversal of
BRAF-vs-RAS biology.

\paragraph{When to worry vs when it is fine.} A practical guide:
(i) gene-level biomarker panel for clinical assay --- watch gene-level
DIAL on the training space; do not deploy under a flip; (ii) pathway
enrichment --- safe to proceed if every pathway has DIAL $\leq 0.1$, as
here; (iii) deep-learning on raw expression --- re-measure
embedding-space DIAL after encoding.

% =====================================================================
\section{DE biomarker recomputation with DESeq2}
\label{sec:S5}

\paragraph{Rationale.} The upstream biomarker list used a Welch
$t$-test on log2-TPM with BH-FDR. Reviewers prefer the counts-based
NB-GLM in DESeq2 \citep{love2014deseq2, muzellec2023pydeseq2}, which
provides size-factor normalisation, empirical-Bayes dispersion
shrinkage, and a Wald test on moderated log2 fold-changes --- all
more robust for low-count genes than a $t$-test on TPM. We reran the
BRAF-vs-RAS contrast with \texttt{pydeseq2} v0.5.4.

\paragraph{Data and design.} TCGA-THCA $n=351$ (293 BRAF, 58 RAS).
Design: \texttt{\textasciitilde subtype}. Cut:
$\mathrm{FDR}_{\mathrm{BH}} < 0.05$, $|\log_2\mathrm{FC}| > 1$. Raw
STAR counts were not on disk; integer pseudo-counts were derived from
log2-TPM via $\mathrm{round}((2^x - 1) \times 50)$ assuming a nominal
$50$M library. Size-factor normalisation absorbs scaling, but gene-level
variance is understated relative to true counts; results are a
reviewer-oriented sanity check rather than a definitive requantification.

\paragraph{Results.}
\begin{table}[htbp]
\centering
\small
\begin{tabular}{lr}\toprule
Metric & Value \\ \midrule
Old significant genes (prior list) & 2{,}773 \\
New significant genes (DESeq2) & 6{,}283 \\
Intersection & 888 \\
Old only & 1{,}885 \\
New only & 5{,}395 \\
Jaccard & 0.109 \\
\bottomrule\end{tabular}
\caption{DE biomarker recomputation: Welch $t$ vs DESeq2.}
\label{tab:s5}
\end{table}

All eight druggable targets from the prior biomarker slate
(TACSTD2, TMPRSS4, PLEKHA6, CYP1B1, LDLR, GABRB2, B3GNT3, PTPRE) are
retained under DESeq2 with $p_{\mathrm{adj}}$ ranging from
$1.9 \times 10^{-222}$ (TACSTD2) to $5.2 \times 10^{-83}$ (LDLR) and
$\log_2\mathrm{FC}$ from $+2.58$ to $+5.29$. No target needs to be
reprioritised on the basis of this reanalysis.

% =====================================================================
\section{Quantum classifier DIAL comparison}
\label{sec:S6}

\paragraph{Goal.} Test whether the v5.1 THCA flip is classifier-paradigm
specific. The primary analysis used linear and tree ensembles only.
We add a third paradigm --- quantum kernel / variational
\citep{havlicek2019qsvm, schuld2019qml, huang2021quantum} --- and
re-measure DIAL on the same five cancers under LODO. If the THCA flip
is task-intrinsic (a cohort-pairing property of the data) it should
reproduce under quantum classifiers; if it is a peculiarity of linear
models it should not.

\paragraph{Methods.} Per-cancer, downsampled to $n \leq 80$ (stratified
by $Y \times B$), PCA(8 components) to project to an 8-qubit Hilbert
space. Three paradigms: classical SVC-RBF (baseline); QSVC (Qiskit-ML
quantum kernel with ZZFeatureMap, reps=2); VQC (PennyLane variational
classifier, 3 strongly entangling layers, 60 training iterations).
DIAL was computed identically to the primary analysis.

\paragraph{Results.}
\begin{table}[htbp]
\centering
\small
\begin{tabular}{llrrrr}\toprule
Cancer & Family & AUC$_{\mathrm{pre}}$ & AUC$_{\mathrm{post}}$ & DIAL & Interpretation \\ \midrule
THCA & QSVC     & 0.549 & 0.599 & 0.000 & no\_signal (underfit) \\
THCA & VQC      & 0.673 & 0.114 & \textbf{0.386} & \textbf{batch\_entangled} \\
THCA & SVC-RBF  & 1.000 & 0.071 & \textbf{0.429} & \textbf{batch\_entangled} \\
SKCM & QSVC     & 0.476 & 0.782 & 0.000 & true\_biology \\
SKCM & VQC      & 0.505 & 0.434 & 0.066 & no\_signal \\
SKCM & SVC-RBF  & 0.421 & 0.526 & 0.000 & no\_signal \\
LGG  & QSVC     & 0.184 & 0.183 & 0.317 & artefact (sign-inverted) \\
LGG  & VQC      & 0.929 & 0.920 & 0.000 & true\_biology \\
LGG  & SVC-RBF  & 0.998 & 1.000 & 0.000 & true\_biology \\
LUAD & SVC-RBF  & 0.774 & 0.847 & 0.000 & true\_biology \\
COAD & SVC-RBF  & 0.714 & 0.789 & 0.000 & true\_biology \\
\bottomrule\end{tabular}
\caption{DIAL under QSVC, VQC, and SVC-RBF. THCA flips under VQC and
SVC-RBF with magnitudes comparable to the classical LogReg/GBM flips.}
\label{tab:s6}
\end{table}

THCA under VQC flips from $\mathrm{AUC}_{\mathrm{pre}}=0.673$ to
$\mathrm{AUC}_{\mathrm{post}}=0.114$ (DIAL $0.386$), matching the
magnitude of the classical flips. SVC-RBF agrees (DIAL $0.429$). The
THCA-specific direction inversion therefore extends from linear +
tree to variational-quantum paradigms; the reviewer objection "DIAL
is a linear-model weakness" is answered negatively. QSVC on THCA
underfits (AUC $\approx 0.55$) --- an underfit, not a flip. LGG/QSVC
shows a sign-encoding artefact (AUC $\approx 0.18$ pre and post) but
VQC and SVC-RBF on the same LGG folds both reach AUC $\geq 0.92$, so
this is classifier-specific, not a batch flip.

% =====================================================================
\section{Han-lab-inspired methodology}
\label{sec:S7}

Three Han-lab-style analyses strengthen the thyroid-specific claim:
(S7.1) FastRNA-style per-cohort centering \citep{lee2022fastrna},
(S7.2) a BUHMBOX-inspired concept check for hidden sub-group
heterogeneity \citep{han2016buhmbox}, and (S7.3) linear-mixed-model
biomarker correction in the style of \citep{sul2018population}. All
use the harmonised THCA cohort (n=392, BRAF=321/RAS=71;
TCGA-THCA=351/GSE27155=41; top-3000 variance genes).

\subsection{FastRNA-style cohort centering}
\label{sec:S7.1}
FastRNA subtracts the per-cohort mean from every sample --- variance
is preserved, global location is removed, no empirical-Bayes step is
required. Applied in place of ComBat and followed by the v5.1 LODO
DIAL pipeline:

\begin{table}[htbp]
\centering
\small
\begin{tabular}{lrrrr}\toprule
Classifier & AUC$_{\mathrm{pre}}$ & AUC$_{\mathrm{post,centered}}$ & DIAL$_{\mathrm{centered}}$ & DIAL$_{\mathrm{ComBat}}$ \\ \midrule
LogReg-$\ell_2$      & 0.994 & 0.993 & 0.000 & 0.494 \\
LogReg-elasticnet    & 0.994 & 0.993 & 0.000 & 0.492 \\
RandomForest         & 0.996 & 0.950 & 0.000 & 0.323 \\
GradientBoosting     & 0.750 & 0.745 & 0.000 & 0.334 \\
XGBoost              & 0.781 & 0.884 & 0.000 & 0.013 \\
\bottomrule\end{tabular}
\caption{FastRNA-style cohort centering eliminates the THCA flip.}
\label{tab:s7a}
\end{table}

ComBat LODO gave mean DIAL $=0.331$ on THCA with two linear models at
$\approx 0.49$ (essentially perfect label polarity flips). FastRNA
centering yields DIAL $= 0.000$ for every classifier. The THCA
cross-cohort signal is recoverable once the (highly unbalanced)
per-cohort mean is removed.

\subsection{BUHMBOX concept check}
\label{sec:S7.2}
BUHMBOX \citep{han2016buhmbox} detects whether shared heritability
between two case groups reflects shared aetiology or hidden sub-group
admixture. Our check is \emph{conceptual}: for each subtype, we fit
PCA on the pooled class-specific expression matrix and tested whether
PC1 scores from TCGA-THCA and GSE27155 follow the same distribution
(two-sample Kolmogorov--Smirnov). We observe KS $= 1.00$ with
$p = 1.3 \times 10^{-40}$ for BRAF ($n = 28$ vs $293$) and
$p = 3.4 \times 10^{-14}$ for RAS ($n = 13$ vs $58$); both flagged
\texttt{hidden\_heterogeneity}. PC1 aligns perfectly with cohort
identity even after restricting to a single subtype.

\subsection{LMM biomarker correction}
\label{sec:S7.3}
We fit per-gene
$\texttt{expression} \sim \texttt{subtype} + (1 \mid \texttt{cohort})$
(\texttt{statsmodels.MixedLM.from\_formula}, REML, L-BFGS) and
compared to naive OLS $\texttt{expression} \sim \texttt{subtype}$.
Both BH-FDR corrected across $3{,}000$ top-variance genes. At
FDR$<0.05$ the naive test flagged $2{,}152$ genes; LMM flagged
$2{,}201$. The lists overlap on $1{,}661$ genes; $491$ naive hits
were lost with LMM and $540$ were gained.

\begin{table}[htbp]
\centering
\small
\begin{tabular}{lrrrrl}\toprule
Status & $n$ & median $|\beta|$ & median $p_{\mathrm{naive}}$ & median $p_{\mathrm{LMM}}$ & reading \\ \midrule
same               & 1{,}969 & 0.559 & --- & --- & strong-effect robust core \\
gained with LMM    & 540     & 0.265 & $3.1 \times 10^{-1}$ & $3.6 \times 10^{-5}$ & rescued biology \\
lost with LMM      & 491     & 0.040 & $2.0 \times 10^{-2}$ & $4.2 \times 10^{-1}$ & tiny-effect OLS artefacts \\
\bottomrule\end{tabular}
\caption{LMM versus naive OLS biomarker lists on the 3{,}000-gene
universe. LMM \emph{rescues} more genes than it drops; dropped genes
are tiny-effect cohort-confounded artefacts.}
\label{tab:s7c}
\end{table}

$99.4\%$ of the prior 2{,}773-gene biomarker slate that is testable
by LMM retains significance under the cohort random-effect correction;
all eight druggable targets remain LMM-significant.

% =====================================================================
\section{Direction-invariance lemma (full proof)}
\label{sec:S8}

\paragraph{Setup.} Work with standardised expression features
$X \in \mathbb{R}^{n \times p}$, binary labels
$Y \in \{0,1\}^n$, and multi-valued batch indicator
$B \in \{1,\ldots,k\}^n$. Let $\mu_Y := \mathbb{E}[X \mid Y = 1] -
\mathbb{E}[X \mid Y = 0] \in \mathbb{R}^p$ be the population
biological direction and $\Delta_B := \{\mu_{B=b} - \mu_{B=b'} :
b \neq b'\}$ the set of pair-wise batch-mean differences, with
$\mathcal{V}_B = \mathrm{span}(\Delta_B)$.

A \emph{linear subtype-preserving batch-correction operator} is any
affine map $P : \mathbb{R}^p \to \mathbb{R}^p$ whose sample-level
action $P(X_i) = X_i - \hat\gamma_{B_i}$ satisfies (i) the
post-correction mean within each batch is identical across batches and
(ii) the within-label means $\mathbb{E}[P(X) \mid Y = c]$ preserve the
pre-correction within-label means up to a label-independent constant.
Constraints (i)--(ii) characterise the subtype-preserving ComBat
family \citep{johnson2007combat} and its population limit.

\begin{lemma}[First-moment annihilation]
\label{lem:main}
Let $P$ satisfy (i)--(ii) and assume $\mu_Y \in \mathcal{V}_B$. Then
\[
  \mathbb{E}[P(X) \mid Y = 1] - \mathbb{E}[P(X) \mid Y = 0] = 0.
\]
\end{lemma}

\paragraph{Proof.} Constraint (i) implies, for every pair of batches
$b \neq b'$,
$\mathbb{E}[P(X) \mid B = b] = \mathbb{E}[P(X) \mid B = b']$. The
batch-mean-difference vectors therefore lie in the kernel of $P$;
equivalently, the restriction $P|_{\mathcal{V}_B}$ is the zero map.
Writing $P = \Pi_{\mathcal{V}_B^\perp} + r$ where $r$ is the
label-independent constant guaranteed by (ii), we have
$P(\mu_Y) = \Pi_{\mathcal{V}_B^\perp}(\mu_Y) + 0 = 0$ under the
hypothesis $\mu_Y \in \mathcal{V}_B$, and taking expectations
conditional on $Y$ yields the stated equality. $\square$

\begin{corollary}[Residual predictive content]
\label{cor:main}
Under the hypotheses of Lemma~\ref{lem:main}, the only expression
direction on which the two classes can be separated in $P(X)$ is
determined by the within-label second-moment residual
\[
  \Sigma^\perp := \Pi_{\mathcal{V}_B^\perp}(\Sigma_1 - \Sigma_0)
  \Pi_{\mathcal{V}_B^\perp},
\]
$\Sigma_c = \mathrm{Cov}(X \mid Y = c)$. Any linear classifier
$f_{\mathrm{post}}$ obtained by empirical risk minimisation with a
strictly convex regulariser on $P(X)$ has its decision direction
aligned with the leading eigenvector of $\Sigma^\perp$. Hence (i)
$\mathrm{auc\_flip}(f_{\mathrm{post}})$ may remain bounded away from
$0.5$ whenever $\Sigma^\perp$ is non-degenerate; (ii) the sign of the
resulting discriminant is determined by spectral convention on
$\Sigma^\perp$, not by $\mu_Y$; and (iii)
$\mathrm{DIAL}(f_{\mathrm{post}}) > 0$ reports the batch-confounded
predictive content that survives with its biological sign lost.
\end{corollary}

\paragraph{Sketch.} Lemma~\ref{lem:main} annihilates the
first-moment label signal, so any linear separator on $P(X)$ must be
recovered from the second-moment structure projected onto
$\mathcal{V}_B^\perp$. Empirical risk minimisation over linear
scoring rules with a strictly convex regulariser (the assumption
$f_{\mathrm{post}}$ satisfies) returns a unique minimiser aligned
with the leading eigenvector of the projected within-label covariance
difference; the sign of that eigenvector is determined by spectral
convention. See \citep{leek2010sva, nygaard2016batch,
korsunsky2019harmony} for empirical over-correction geometry that
matches this residual-second-moment regime. $\square$

\paragraph{What the lemma does \emph{not} prove.} We are careful
not to overclaim. The lemma states that the first-moment label
signal is exactly annihilated when $\mu_Y \in \mathcal{V}_B$. It
does not state, and we do not prove, a probabilistic AUC flip of the
form $\mathrm{AUC}(f_{\mathrm{post}}) \to 1 -
\mathrm{AUC}(f_{\mathrm{pre}})$ with probability one under some
generating distribution. Such a statement would require the
held-out-cohort data to be sampled from a distribution whose
spectral gap between $\Sigma_1^\perp$ and $\Sigma_0^\perp$ is
bounded below --- a condition neither TCGA nor GEO are guaranteed to
satisfy, and whose verification would itself require repeated-cohort
simulation. The empirical five-cancer specificity pattern we report
in \S\ref{sec:S3} is the substantive claim of the paper; the
lemma+corollary pair provide the geometric reason why DIAL is the
right thing to measure, not a distributional guarantee.

\paragraph{Remark on identifiability.} When $\mu_Y \in \mathcal{V}_B$,
a separate multi-class classifier trained to predict $B$ from $P(X)$
will attain AUC close to chance because the batch means have been
removed. The pathological regime is therefore the pair (high DIAL,
low batch identifiability), which the primary analysis reports where
the auxiliary classifier is fitable (three-or-more-cohort collections;
not the present two-cohort LODO).

% =====================================================================
\section{Code and data availability}
\label{sec:S9}

\paragraph{Code.} All pipeline code is written in Python 3.11 and is
available in the project repository under
\texttt{notebooks\_or\_scripts/v5p1\_*.py} and
\texttt{notebooks\_or\_scripts/v8\_*.py}. The harmonisation, ComBat,
ComBat-seq, LODO, and meta-analysis steps are each stand-alone
scripts with deterministic seeds.

\paragraph{Data.} Inputs are downloaded at runtime from the GDC data
portal and the NCBI GEO FTP server; no synthetic substitutes are
used. TCGA mutation annotations were obtained from the MC3 public
MAF. Per-cohort harmonised matrices and per-classifier LODO results
are exported as TSV under \texttt{results/v5/} and
\texttt{results/v8\_statgen/}.

\paragraph{Dashboard.} An interactive Plotly-based dashboard hosts
every figure at gene and pathway level and is available at
\url{http://40.82.129.113:8012/}.

\paragraph{Reproducibility.} A \texttt{Dockerfile} and
\texttt{docker-compose.yml} under the project root reproduce the
full pipeline end-to-end from public sources. A frozen pip
\texttt{requirements.txt} is shipped with the image.

% =====================================================================
\bibliographystyle{plainnat}
\bibliography{../references}

\end{document}
